\subsection{Метод с постоянным шагом}

\begin{theorem}\label{thm:const-step}
Пусть $\alpha_k = \alpha = \mathrm{const}$, $0 < \alpha < \dfrac{2}{M}$. Тогда для
скорости сходимости справедлива оценка
\[
  \norm{x_{k+1} - x_*} \ \le\ q\, \norm{x_k - x_*} ,
  \qquad
  q = \max\big\{\, \abs{1 - \alpha m},\ \abs{1 - \alpha M} \,\big\} ,
\]
причём
\[
  q_{\min} = \frac{M - m}{M + m}
  \quad \text{при} \quad
  \alpha = \frac{2}{M + m} ,
\]
и, кроме того,
\[
  f(x_k) - f(x_*) \ \le\ C\, q^{2k} , \qquad C < +\infty .
\]
\end{theorem}

\begin{proof}
Так как $\nabla f(x_*) = 0$, по формуле конечных приращений
\[
  \nabla f(x_k) - \nabla f(x_*) = H(\bar{x}_k)\,(x_k - x_*) ,
  \qquad \bar{x}_k = x_k + \theta\,(x_* - x_k), \quad \theta \in (0,1) .
\]
Поэтому
\begin{align*}
  \norm{x_{k+1} - x_*}
  &= \norm{x_k - \alpha \nabla f(x_k) - x_*}
   = \norm{x_k - x_* - \alpha\big(\nabla f(x_k) - \nabla f(x_*)\big)} \\
  &= \norm{x_k - x_* - \alpha H(\bar{x}_k)(x_k - x_*)}
   = \norm{\big(E - \alpha H(\bar{x}_k)\big)(x_k - x_*)} \\
  &\le \norm{E - \alpha H(\bar{x}_k)} \cdot \norm{x_k - x_*} .
\end{align*}
Обозначив $q = \norm{E - \alpha H(\bar{x}_k)}$, получаем
$\norm{x_{k+1} - x_*} \le q \norm{x_k - x_*}$.

Матрица $H(\bar{x}_k)$ симметрична, и её собственные числа лежат в $[m, M]$.
Тогда $E - \alpha H(\bar{x}_k)$ симметрична с собственными числами
$1 - \alpha \mu$, $\mu \in [m, M]$, и её спектральная норма равна
\[
  \norm{E - \alpha H(\bar{x}_k)}
  = \max_{\mu \in [m,M]} \abs{1 - \alpha \mu}
  = \max\big\{\, \abs{1 - \alpha m},\ \abs{1 - \alpha M} \,\big\} ,
\]
поскольку модуль линейной функции достигает максимума на конце отрезка. Эта
оценка не зависит от $\bar{x}_k$, то есть одна и та же на всех итерациях.

Минимум по $\alpha$ достигается там, где две ветви сравниваются:
\[
  1 - \alpha m = \alpha M - 1
  \qquad\Longrightarrow\qquad
  \alpha = \frac{2}{m + M} ,
  \qquad
  q_{\min} = 1 - \frac{2m}{m+M} = \frac{M - m}{M + m} .
\]

Наконец, из правой части~\eqref{eq:rate-2}
\[
  f(x_k) - f(x_*) \ \le\ \frac{M}{2}\norm{x_k - x_*}^2
  \ \le\ \frac{M}{2}\, q^{2k} \norm{x_0 - x_*}^2 = C\, q^{2k} ,
  \qquad
  C = \frac{M}{2}\norm{x_0 - x_*}^2 .
  \qedhere
\]
\end{proof}

\inputfig{figures/fig-q-alpha}

\begin{idea}[title={Два места в доказательстве, которые стоит разобрать}]
\textbf{1. Переход $\norm{A\,y} \le \norm{A}\cdot\norm{y}$.} Это определение
операторной (подчинённой) нормы матрицы:
$\norm{A} = \sup_{y \ne 0} \norm{Ay} / \norm{y}$, откуда
$\norm{Ay} \le \norm{A}\norm{y}$ для всех $y$ автоматически. Для евклидовой
нормы вектора подчинённая норма матрицы — спектральная, а для симметричной
матрицы она равна максимуму модулей собственных чисел. Именно поэтому в
следующей строке можно перейти от нормы матрицы к собственным числам гессиана.

\medskip
\textbf{2. Условие $0 < \alpha < 2/M$.} Метод сходится тогда и только тогда,
когда $q < 1$, то есть когда \emph{обе} величины $\abs{1 - \alpha m}$ и
$\abs{1 - \alpha M}$ меньше единицы. Так как $m \le M$, ограничение даёт именно
большая из них: $\abs{1 - \alpha M} < 1 \iff 0 < \alpha < 2/M$. При
$\alpha \ge 2/M$ шаг «перелетает» минимум и метод расходится.

\medskip
\textbf{Что показывает график.} Ветвь $1 - \alpha m$ убывает и обнуляется в
$\alpha = 1/m$; ветвь $\abs{1 - \alpha M}$ обнуляется раньше, в $\alpha = 1/M$,
и дальше растёт. Максимум из двух сначала идёт по первой ветви, после точки
пересечения — по второй, поэтому минимум максимума достигается ровно в точке
пересечения $\alpha = 2/(m+M)$.
\end{idea}

\begin{idea}
Сравним три метода. Все оценки выражаются через $t = m/M \in (0,1]$:
\begin{gather*}
  \text{дробление и наискорейший спуск:} \quad q = 1 - \tfrac{t}{2}(1+t) , \\
  \text{постоянный шаг:} \quad q^2 = \Big(\frac{1-t}{1+t}\Big)^2 .
\end{gather*}
Обе величины стремятся к единице при $t \to 0$, то есть все три метода одинаково
страдают от плохой обусловленности. Разница в цене итерации: наискорейший спуск
требует решения одномерной задачи, дробление — нескольких пробных вычислений
$f$, постоянный шаг — ни того, ни другого, но требует знания $m$ и $M$, которые
на практике редко известны. Радикально ситуацию меняет только переход к методам
второго порядка.
\end{idea}
